\chapter{RESULTS}

This chapter explains the result of the implemented Inventory Management System. The result is based on the current codebase and the pages that are already implemented in the web application. The detailed development method and calculation logic were explained in Chapter 3, so this chapter focuses on what the completed system can do.

\section{Completed System Overview}
The final system is a full-stack web application for an SME trading business. It supports the flow from product setup to purchasing, selling, billing, supplier payment, credit note adjustment, inventory checking, dashboard reporting, and limited AI-assisted questions.

The main result is that the system connects business documents together instead of keeping them as separate records. A product can appear in purchases, sales, quotations, inventory reports, and product history. A sale can later become part of a billing note or a credit note. A purchase can later become part of a payment batch. This connection makes the system more useful than a simple product list.

\section{Implemented Modules}
\begin{table}[H]
	\centering
	\small
	\renewcommand{\arraystretch}{1.18}
	\setlength{\tabcolsep}{4pt}
	\caption{Implemented Workspace and Transaction Modules}
	\label{tab:implemented-system-modules}
	\begin{tabularx}{\textwidth}{|p{3cm}|X|p{2.1cm}|}
		\hline
		\textbf{Module} & \textbf{Implemented Result} & \textbf{Status} \\ \hline
		Dashboard & Shows stock, finance, cashflow, product movement, and order coverage summaries. & Completed \\ \hline
		Inventory & Shows backend-calculated stock rows, reorder need, supplier options, product details, and quick purchase-order preparation. & Completed \\ \hline
		AI Chat & Answers supported read-only questions about stock, fulfillment, partner summaries, AR/AP exceptions, and document references. & Completed \\ \hline
		Purchases & Supports purchase records, line items, receiving status, documents, discounts, VAT totals, payable total, and payment batch links. & Completed \\ \hline
		Payment Batches & Groups eligible purchases for supplier payment and tracks paid or unpaid lines. & Completed \\ \hline
		Quotations & Supports quotation records, line items, supplier options, validity settings, VAT totals, and conversion to purchase or sale flow. & Completed \\ \hline
		Sales & Supports sales records, stock validation, stock allocation, delivery status, documents, VAT totals, billing links, and credit links. & Completed \\ \hline
	\end{tabularx}
\end{table}

\begin{table}[H]
	\centering
	\small
	\renewcommand{\arraystretch}{1.18}
	\setlength{\tabcolsep}{4pt}
	\caption{Implemented Finance and Master Data Modules}
	\label{tab:implemented-record-modules}
	\begin{tabularx}{\textwidth}{|p{3cm}|X|p{2.1cm}|}
		\hline
		\textbf{Module} & \textbf{Implemented Result} & \textbf{Status} \\ \hline
		Billing Notes & Groups eligible sales for customer receivable tracking and tracks received or outstanding lines. & Completed \\ \hline
		Credit Notes & Creates credit notes from eligible cancelled or returned sale items. & Completed \\ \hline
		Products & Supports product CRUD, SKUs, pictures, unit conversions, reorder levels, active status, supplier links, and product history. & Completed \\ \hline
		Categories & Supports category records with parent-child hierarchy. & Completed \\ \hline
		Suppliers & Supports supplier profile, contact, branch, shipping address, procurement contact, and payment term data. & Completed \\ \hline
		Customers & Supports customer profile, contact, branch, shipping address, and billing term data. & Completed \\ \hline
	\end{tabularx}
\end{table}

\section{Screenshot Placeholders}
The following figures are reserved for the final screenshots. They should be replaced with real screenshots from the completed application after the final demo data is approved.

\begin{figure}[H]
\centering
\fbox{\parbox{0.8\textwidth}{TODO: Insert screenshot of the dashboard page showing stock and finance summary cards.}}
\caption{Dashboard page}
\label{fig:dashboard-page}
\end{figure}

\begin{figure}[H]
\centering
\fbox{\parbox{0.8\textwidth}{TODO: Insert screenshot of the inventory page showing stock rows, filters, and reorder information.}}
\caption{Inventory page}
\label{fig:inventory-page}
\end{figure}

\begin{figure}[H]
\centering
\fbox{\parbox{0.8\textwidth}{TODO: Insert screenshot of the product form showing product details, unit conversions, and product picture support.}}
\caption{Product management form}
\label{fig:product-management-form}
\end{figure}

\begin{figure}[H]
\centering
\fbox{\parbox{0.8\textwidth}{TODO: Insert screenshot of the purchase workflow showing purchase items and receiving status.}}
\caption{Purchase workflow}
\label{fig:purchase-workflow}
\end{figure}

\begin{figure}[H]
\centering
\fbox{\parbox{0.8\textwidth}{TODO: Insert screenshot of the sales workflow showing delivery status and stock allocation.}}
\caption{Sales workflow}
\label{fig:sales-workflow}
\end{figure}

\begin{figure}[H]
\centering
\fbox{\parbox{0.8\textwidth}{TODO: Insert screenshot of the quotation page showing quotation items and supplier options.}}
\caption{Quotation workflow}
\label{fig:quotation-workflow}
\end{figure}

\begin{figure}[H]
\centering
\fbox{\parbox{0.8\textwidth}{TODO: Insert screenshot of the billing note page showing eligible sales or billing note detail.}}
\caption{Billing note workflow}
\label{fig:billing-note-workflow}
\end{figure}

\begin{figure}[H]
\centering
\fbox{\parbox{0.8\textwidth}{TODO: Insert screenshot of the payment batch page showing eligible purchases or payment batch detail.}}
\caption{Payment batch workflow}
\label{fig:payment-batch-workflow}
\end{figure}

\begin{figure}[H]
\centering
\fbox{\parbox{0.8\textwidth}{TODO: Insert screenshot of a supported read-only assistant response, such as stock status or reference lookup.}}
\caption{AI inventory assistant}
\label{fig:ai-inventory-assistant}
\end{figure}

\section{Results by Workflow}
\subsection{Master Data Results}
The master data pages were completed for products, categories, suppliers, and customers. Products can store SKUs, previous SKUs, names, categories, stock base units, default purchase and sales units, reorder levels, pictures, supplier links, and active status. Categories support hierarchy through parent categories. Suppliers and customers store contact, branch, shipping address, taxpayer, remark, and payment term information.

This result is important because the transaction pages depend on these records. Purchases need suppliers and products. Sales need customers and products. Quotations can use products, suppliers, and customers. The master data pages therefore work as the foundation of the system.

\subsection{Purchasing and Stock Results}
The purchase workflow was completed with purchase headers, purchase line items, item receiving status, expected delivery date, received date, uploaded documents, VAT fields, discounts, grand total, and payable total. A received purchase item increases available stock, while pending and cancelled purchase items do not.

The stock result is calculated by the backend. The inventory and dashboard pages show current stock, pending purchase units, delayed purchase units, pending sales units, recommended restock, stock value, and supplier options. This makes the inventory page useful for daily checking because the user does not need to manually combine purchase and sales records.

\subsection{Sales Results}
The sales workflow was completed with sale headers, customer information, customer purchase order reference, sale line items, delivery statuses, uploaded documents, VAT fields, discounts, billing note links, credit note links, and source quotation links.

The most important result in the sales module is server-side stock validation. A sale can stay as a draft even if there is not enough stock, but it cannot commit unavailable stock when the sale item becomes packed, shipped, or delivered. When stock is committed, the backend creates allocation records from received purchase layers. This gives better traceability between stock received from suppliers and stock sold to customers.

\subsection{Quotation Results}
The quotation workflow was completed with quotation date, validity settings, customer and supplier references, shipping date, payment terms, VAT mode, total values, line items, and supplier options for each line. Quotation lines are stored in relational tables, not as one large JSON field. The frontend can also start purchase or sale creation from a quotation.

\subsection{Finance Results}
The finance workflow was completed through billing notes, payment batches, and credit notes. Billing notes group eligible sales for customer receivables. Payment batches group eligible purchases for supplier payables. Credit notes are created from eligible cancelled or returned sale items.

The backend eligibility endpoints are one of the useful results of this implementation. They reduce the chance of selecting the same sale into multiple active billing notes, selecting the same purchase into multiple active payment batches, or crediting the same returned/cancelled sale item more than once.

\subsection{Dashboard and Inventory Results}
The dashboard and inventory pages show the business status without loading every transaction into the frontend. The backend prepares stock report rows, finance summaries, cashflow buckets, product movement, popular products, and order coverage data. This result makes the interface faster and easier to use because the frontend receives summary data that is already prepared for display.

\subsection{AI Assistant Results}
The AI assistant was implemented as a limited read-only assistant. It can answer supported questions about low stock, restock scope, net position, order coverage, customer summary, supplier summary, overdue or exception issues, and line items for a known reference. The tests also show that questions outside these supported workflows are rejected as outside the current assistant scope.

\section{Feature Evidence Summary}
\begin{table}[H]
	\centering
	\small
	\renewcommand{\arraystretch}{1.25}
	\setlength{\tabcolsep}{4pt}
	\caption{Evidence of Implemented Features}
	\label{tab:feature-evidence-summary}
	\begin{tabularx}{\textwidth}{|p{3.7cm}|X|}
		\hline
		\textbf{Feature} & \textbf{Evidence in the Implemented System} \\ \hline
		REST API communication & The frontend API client calls endpoints for dashboard, lookups, products, purchases, sales, quotations, billing notes, payment batches, credit notes, authentication, and chat. \\ \hline
		JWT authentication support & The backend provides login, refresh, and profile endpoints, and the frontend stores and refreshes access tokens. \\ \hline
		Backend stock validation & The sale workflow calls backend validation before stock-deducting sale statuses are committed. \\ \hline
		FIFO and manual allocation & The backend creates sale allocations from received purchase layers and validates manual layer selections. \\ \hline
		Finance eligibility & Backend endpoints return only eligible sales for billing notes, eligible purchases for payment batches, and eligible sale lines for credit notes. \\ \hline
		Normalized quotation items & Quotation lines and supplier options are stored in relational tables while the API still returns a simple \texttt{items} array to the frontend. \\ \hline
		Backend dashboard summaries & The dashboard and inventory pages use backend-prepared summary data instead of calculating all operational metrics from full transaction lists in the browser. \\ \hline
		Assistant scope control & The chat service answers only supported read-only workflows and rejects unsupported requests as outside scope. \\ \hline
	\end{tabularx}
\end{table}

\section{Key Findings}
The first key finding is that inventory stock should be derived from transactions, not typed manually into the product record. In this system, stock depends on received purchase items and committed sales. This makes stock more reliable because it follows the actual business documents.

The second key finding is that finance documents need backend eligibility rules. Billing notes, payment batches, and credit notes are easy to duplicate by mistake if the frontend is the only place checking them. Moving the rules to the backend makes the workflow safer.

The third key finding is that historical snapshot fields are useful. Even when products, suppliers, and customers are stored as master data, old transactions still need to keep names, SKUs, prices, and totals so that the document remains understandable later.

\section{Discussion of Objective Achievement}
The first objective was to record purchase and sales transactions with process status and attachment support. This objective is achieved because purchases and sales both support line items, statuses, document uploads, and transaction detail views.

The second objective was to provide inventory stock information and calculation support. This objective is achieved through backend stock calculations, inventory page stock rows, dashboard summaries, product history, stock layers, reorder information, and recommended restock values.

The third objective was to integrate an AI chatbot connected to backend data. This objective is achieved within the current implemented scope. The assistant can answer supported read-only operational questions from backend data, but it is not a general-purpose chatbot and it does not change system records.

\section{User Survey Evaluation}
At the current stage, the user survey evaluation has not been conducted because no users have evaluated the software yet. This section is kept as a placeholder for the final blackbook after real users or classmates have tested the system.

\begin{table}[H]
	\centering
	\small
	\renewcommand{\arraystretch}{1.25}
	\setlength{\tabcolsep}{4pt}
	\caption{User Survey Evaluation Placeholder}
	\label{tab:user-survey-placeholder}
	\begin{tabularx}{\textwidth}{|p{4cm}|X|}
		\hline
		\textbf{Evaluation Item} & \textbf{Result} \\ \hline
		Number of participants & TODO: Insert number of users who evaluated the system. \\ \hline
		User group & TODO: Insert participant type, such as students, SME staff, or project reviewers. \\ \hline
		Evaluation method & TODO: Insert survey method, rating scale, and evaluation questions. \\ \hline
		Average score & TODO: Insert average usability or satisfaction score after the survey is completed. \\ \hline
		User comments & TODO: Insert summarized feedback from users. \\ \hline
	\end{tabularx}
\end{table}

% TODO: Replace this placeholder with real user survey results after users evaluate the software.

\section{Limitations Found During Development}
The system is functional, but some limitations remain:

\begin{itemize}
	\item Final screenshots still need to be inserted into the report.
	\item User survey results are currently placeholders because no users have evaluated the software yet.
	\item Backend-wide authentication is configurable, so the final deployment setting must be checked before production use.
	\item The frontend still keeps mock-data fallback behavior for demonstration, so final evaluation should use backend-connected data.
	\item Large-scale production performance testing with many users and a very large database was not completed in the current codebase.
\end{itemize}

% TODO: Insert final experiment results, user survey results, performance notes, and approved screenshots if they become available.
